\section{Proofs and Mathematical Audit}
\label{app:proofs}

For any positive-semidefinite matrix $\bH$, write
$\norm{\bm{v}}_{\bH}^2=\bm{v}^\top\bH\bm{v}$.

\subsection{Proof of Lemma~\ref{lem:local-forgetting}}

\begin{proof}
Let $\bm{d}=\bx_t-\bx_s$. Taylor's theorem with integral remainder gives
\begin{align}
\cL_s(\bx_s+\bm{d})
&=\cL_s(\bx_s)+\nabla\cL_s(\bx_s)^\top\bm{d}
+\frac12\bm{d}^\top\nabla^2\cL_s(\bx_s)\bm{d}+r_{s,t},
\end{align}
where
\begin{equation}
r_{s,t}=\int_0^1(1-\tau)\bm{d}^\top
\left[\nabla^2\cL_s(\bx_s+\tau\bm{d})-\bH_s\right]
\bm{d}\,d\tau.
\end{equation}
Hessian Lipschitzness implies
\begin{align}
|r_{s,t}|
&\leq \rho_s\norm{\bm{d}}_2^3
\int_0^1\tau(1-\tau)\,d\tau
=\frac{\rho_s}{6}\norm{\bm{d}}_2^3.
\end{align}
Finally, $\bg_s^\top\bm{d}=0$ for every $\bm{d}\in\cU$ if and only if
$\bg_s\in\cU^\perp$, equivalently $\Pi_{\cU}\bg_s=\bm{0}$.
\end{proof}

For a nonquadratic old loss, the one-step change has the remainder difference
\begin{equation}
\cF_{s\to t}-\cF_{s\to t-1}
=\bg_s^\top\bh_t
+(\bx_{t-1}-\bx_s)^\top\bH_s\bh_t
+\frac12\bh_t^\top\bH_s\bh_t+\Delta r_{s,t},
\end{equation}
with the direct bound
\begin{equation}
|\Delta r_{s,t}|
\leq\frac{\rho_s}{6}
\left(\norm{\bx_t-\bx_s}_2^3
+\norm{\bx_{t-1}-\bx_s}_2^3\right).
\end{equation}
Consequently, the quadratic increment is informative only when this remainder
is controlled relative to the leading terms.

\subsection{Proof of Sequential Accumulation Results}

\begin{proof}[Proof of Theorem~\ref{thm:sequential-accumulation}]
The telescoping identity
$\bx_t-\bx_s=\sum_{u=s+1}^{t}\bh_u$ substituted into
Eq.~\eqref{eq:quadratic-old-loss} gives
Eq.~\eqref{eq:path-forgetting}. Since
$\bx_t-\bx_s=(\bx_{t-1}-\bx_s)+\bh_t$, expanding the squared seminorm gives
Eq.~\eqref{eq:forgetting-increment}. Taking expectations, applying
Eq.~\eqref{eq:noncancellation}, and summing from $s+1$ to $T$ yields
\begin{align}
\E\cF_{s\to T}
&\geq\sum_{t=s+1}^{T}
\left(\frac12\E\norm{\bh_t}_{\bH_s}^2-\beta_t\right).
\end{align}
The stated special case follows from $\beta_t=0$ and the lower bounds
$\nu_t^2$.
\end{proof}

\begin{proof}[Proof of Corollary~\ref{cor:independent-innovations}]
Let $\bm{S}=\sum_{u=s+1}^{T}\bh_u$. Independence gives
$\E\bm{S}=\sum_u\bm{\mu}_u$ and
$\operatorname{Cov}(\bm{S})=\sum_u\bSigma_u$. For any random vector with a
finite second moment,
\begin{equation}
\E[\bm{S}^\top\bH_s\bm{S}]
=(\E\bm{S})^\top\bH_s(\E\bm{S})
+\tr(\bH_s\operatorname{Cov}(\bm{S})).
\end{equation}
Multiplying by $1/2$ and invoking Eq.~\eqref{eq:path-forgetting} proves the
claim.
\end{proof}

\begin{proof}[Proof of Proposition~\ref{prop:mixture-bridge}]
Completing the square in the mixture risk gives
\begin{equation}
\cL_t(\bx)=\kappa_t'
+\frac12\norm{\bx-\bm{m}_t}_{\bH}^2,
\end{equation}
where $\kappa_t'$ is independent of $\bx$. The constrained minimizer is therefore
$P_{\cU}^{\bH}\bm{m}_t$. Let
$\bm{d}=P_{\cU}^{\bH}(\bm{m}_t-\bm{m}_s)\in\cU$. Because
$\bx_s-\bm{m}_s=-(\bI-P_{\cU}^{\bH})\bm{m}_s$ is $\bH$-orthogonal to
$\cU$,
\begin{align}
\cL_s(\bx_t)-\cL_s(\bx_s)
&=\frac12\norm{(\bx_s-\bm{m}_s)+\bm{d}}_{\bH}^2
-\frac12\norm{\bx_s-\bm{m}_s}_{\bH}^2\\
&=\frac12\norm{\bm{d}}_{\bH}^2.
\end{align}
Finally,
$\bm{m}_t-\bm{m}_s=\bM\sum_{u=s+1}^{t}\Delta\bm{\pi}_u$, which yields the
increment formula.
\end{proof}

\begin{remark}[Cancellation counterexample]
Let $p=1$, $H_s=1$, $h_{s+1}=1$, and $h_{s+2}=-1$. Then
$\cF_{s\to s+1}=1/2$ but $\cF_{s\to s+2}=0$. The cross term at the second
step equals $-1$ and dominates the new quadratic energy $1/2$. Thus no theorem
can infer monotone forgetting solely from sequential ordering.
\end{remark}

\subsection{Proof of the Three-Way Decomposition}

\begin{proof}[Proof of Proposition~\ref{prop:error-decomposition}]
Add and subtract $\overline{\cL}_T(\bX^\star)$ and
$\overline{\cL}_T(\bX_R^\star)$ from the left-hand side. This gives
Eq.~\eqref{eq:three-way-decomposition} algebraically. For nonnegativity of the
conflict term,
\begin{equation}
\min_{\bX}\sum_s w_s\cL_s(\bX)
\geq \sum_s w_s\min_{\bX}\cL_s(\bX).
\end{equation}
The capacity term is nonnegative because $\cC_R$ is contained in the full
update space. The online term is nonnegative because $\bX_R^\star$ globally
minimizes $\overline{\cL}_T$ over $\cC_R$ and
$\widehat{\bX}_T\in\cC_R$.
\end{proof}

\begin{remark}[Conflict survives unlimited rank]
For equal weights and
$\cL_1(x)=\frac12(x-1)^2$, $\cL_2(x)=\frac12(x+1)^2$, the individual minima
have zero loss, while the shared optimum $x^\star=0$ has average loss $1/2$.
The capacity and online terms can both be zero. Hence the retention gap is not,
in general, entirely a capacity-allocation problem.
\end{remark}

\subsection{Proof of the Spectral Capacity Formula}

\begin{proof}[Proof of Theorem~\ref{thm:spectral-capacity}]
Define the invertible change of variables
\begin{equation}
\bY=\bC_{\mathrm{out}}^{1/2}\bX\bC_{\mathrm{in}}^{1/2}.
\end{equation}
Multiplication by invertible matrices preserves rank, so the constrained
capacity problem is equivalent to
\begin{equation}
\min_{\rank(\bY)\leq R}\frac12\norm{\bY-\bY^\star}_F^2.
\end{equation}
The Eckart--Young--Mirsky theorem gives $[\bY^\star]_R$ as a minimizer and
$\frac12\sum_{j>R}\sigma_j^2(\bY^\star)$ as its residual. Applying the inverse
change of variables gives the stated $\bX_R^\star$.
\end{proof}

\begin{proof}[Proof of Corollary~\ref{cor:capacity-lower-bound}]
At least $m-R$ singular values occur in the tail after rank $R$, and each has
squared magnitude at least $\alpha^2$. Substitution into
Eq.~\eqref{eq:spectral-capacity} proves the bound. The sequential statement
additionally assumes that the modes persist rather than cancel or lose energy.
\end{proof}

\begin{proof}[Proof of Corollary~\ref{cor:global-layer-allocation}]
For each layer, apply the invertible whitening transformation from
Theorem~\ref{thm:spectral-capacity}. If layer $\ell$ receives integer rank
$R_\ell$, Eckart--Young--Mirsky gives residual
$\frac12\sum_{j>R_\ell}\sigma_j^2(\bY_\ell^\star)$. Because every retained
singular component consumes one unit of the global budget and removes exactly
one nonnegative squared singular value from the residual, minimizing the sum
under $\sum_\ell R_\ell\leq R$ is achieved by retaining the $R$ largest values
across all layers. Ties may produce multiple optimal allocations.
\end{proof}

\begin{corollary}[Locally optimal rank-constrained consolidation]
\label{cor:merge-consolidation}
Under Eq.~\eqref{eq:kronecker-quadratic}, replacing an unconstrained accumulated
update $\bX^\star$ by Eq.~\eqref{eq:weighted-svd} minimizes historical mixture
loss among all aligned rank-$R$ updates. If
$\sigma_R(\bY^\star)>\sigma_{R+1}(\bY^\star)$, the whitened rank-$R$ subspace is
unique.
\end{corollary}

The corollary is a local statement in a shared parameter basis. It does not
cover independently trained models that require representation alignment, nor
does it control the Taylor remainder after a large merge. Applied to an
arbitrary $\bX\neq\bX^\star$, Eq.~\eqref{eq:weighted-svd} minimizes weighted
reconstruction error to $\bX$, not historical mixture loss itself.

\subsection{Proof of Safe-Space Contraction}

\begin{proof}[Proof of Theorem~\ref{thm:safe-subspace}]
For any ambient displacement $\bm{v}$,
\begin{equation}
\bm{v}^\top\bA_t^{\mathrm{amb}}\bm{v}
=\sum_{s=1}^{t}\omega_s\bm{v}^\top\bK_s\bm{v}.
\end{equation}
Every summand is nonnegative. The sum is zero if and only if every summand is
zero. For a positive-semidefinite matrix,
$\bm{v}^\top\bK_s\bm{v}=0$ if and only if
$\bm{v}\in\ker(\bK_s)$. This proves the ambient kernel identity and nesting.

For the approximate construction,
$\cN_t^{\mathrm{amb},(\tau)}=
\cN_{t-1}^{\mathrm{amb},(\tau)}\cap
\cS_t^{\mathrm{amb}}(\tau)^\perp$, so nesting follows directly. Replacing any
nonzero $\bK_s$ by $c\bK_s$ for $c>0$ scales both its eigenvalues and
$\norm{\bK_s}_{\mathrm{op}}$ by $c$ and therefore leaves
$\cS_s^{\mathrm{amb}}(\tau)$ unchanged. The zero-matrix convention is unchanged
as well.

Finally, for fixed $\bB$,
\begin{equation}
\ba^\top\bA_t^{\bB}\ba
=(\bB\ba)^\top\bA_t^{\mathrm{amb}}(\bB\ba).
\end{equation}
This quantity is zero if and only if
$\bB\ba\in\ker(\bA_t^{\mathrm{amb}})$. Applying the same positive-semidefinite
sum argument to $\bG_s=\bB^\top\bK_s\bB$ proves the remaining identities and
nesting in Eq.~\eqref{eq:coefficient-safe-kernel}.
\end{proof}

The exact remaining coefficient dimension within a fixed dictionary can also
be written as
\begin{equation}
\operatorname{dim}(\cN_t^{\bB})
=m-\rank\begin{bmatrix}
\sqrt{\omega_1}\bG_1^{1/2}\\
\vdots\\
\sqrt{\omega_t}\bG_t^{1/2}
\end{bmatrix}.
\label{eq:safe-dimension}
\end{equation}
Under the explicit general-position condition that the stacked rank equals
$\min\{m,\sum_{s\leq t}\rank(\bG_s)\}$,
\begin{equation}
\operatorname{dim}(\cN_t^{\bB})
=\max\{0,m-\sum_{s\leq t}\rank(\bG_s)\}.
\end{equation}
This dimension-exhaustion rate is conditional on general position; the kernel
nesting itself is unconditional for additive positive-semidefinite terms.

\begin{proof}[Proof of Proposition~\ref{prop:price-useful-update}]
Let $\bz_0=\Pi_{\ker(\bA)}\bz$. If $\bz_0\neq0$, choose
$\ba=\gamma\bz_0/\norm{\bz_0}_2^2$. Then
$\bz^\top\ba=\gamma$ and $\ba^\top\bA\ba=0$, proving the first case.

Otherwise $\bz\perp\ker(\bA)$, hence $\bz\in\range(\bA)$ for symmetric
$\bA$. Generalized Cauchy--Schwarz gives
\begin{equation}
(\bz^\top\ba)^2
\leq(\bz^\top\bA^\dagger\bz)(\ba^\top\bA\ba).
\end{equation}
Any feasible point therefore has cost at least
$\gamma^2/(2\bz^\top\bA^\dagger\bz)$. The displayed $\ba^\star$ in
Eq.~\eqref{eq:price-solution} is feasible and attains equality.
\end{proof}

\subsection{Proof of the Top-$k$ Allocation Rule}

\begin{proof}[Proof of Theorem~\ref{thm:topk-allocation}]
For a fixed support $\cS$, generalized Cauchy--Schwarz yields
\begin{equation}
G(\ba)^2
\leq\left(\sum_{j\in\cS}\bz_j^\top
\widetilde{\bA}_j^{-1}\bz_j\right)
\left(\sum_{j\in\cS}\ba_j^\top
\widetilde{\bA}_j\ba_j\right).
\end{equation}
Thus any update achieving gain $\gamma$ has cost at least
$\gamma^2/(2\sum_{j\in\cS}q_j)$. The coefficient vector in
Eq.~\eqref{eq:topk-solution} attains equality. Minimizing cost over supports of
size at most $k$ is consequently equivalent to maximizing
$\sum_{j\in\cS}q_j$, achieved by the $k$ largest positive scores.
\end{proof}

\begin{corollary}[Robustness to cross-block curvature]
\label{cor:approx-block-allocation}
Let $\bD=\operatorname{blkdiag}(\widetilde{\bA}_1,\ldots,
\widetilde{\bA}_M)$ and suppose the true positive-definite screening curvature
satisfies
\begin{equation}
(1-\delta)\bD\preceq\widetilde{\bA}\preceq(1+\delta)\bD,
\qquad 0\leq\delta<1.
\end{equation}
Let $\widehat{\cS}$ be the top-$k$ support under $\bD$ and $\cS^\star$ the
optimal support under $\widetilde{\bA}$. Write
$Q_{\widetilde{\bA}}(\cS)$ for the minimum true
historical cost at a fixed target gain on support $\cS$. Then
\begin{equation}
Q_{\widetilde{\bA}}(\widehat{\cS})
\leq\frac{1+\delta}{1-\delta}
Q_{\widetilde{\bA}}(\cS^\star).
\end{equation}
\end{corollary}

\begin{proof}
For each support $\cS$, restriction of the spectral sandwich to that support
and inversion imply
\begin{equation}
\frac{1}{1+\delta}\bz_{\cS}^\top\bD_{\cS}^{-1}\bz_{\cS}
\leq
\bz_{\cS}^\top\widetilde{\bA}_{\cS}^{-1}\bz_{\cS}
\leq
\frac{1}{1-\delta}\bz_{\cS}^\top\bD_{\cS}^{-1}\bz_{\cS}.
\end{equation}
The top-$k$ support maximizes the utility surrogate under $\bD$. Combining the
two inequalities for $\widehat{\cS}$ and $\cS^\star$, then using minimum cost
$\gamma^2/(2\,\text{utility})$, gives the factor.
\end{proof}

\begin{remark}[Top-$k$ can fail with strong cross-block terms]
Take
\begin{equation}
\widetilde{\bA}=\begin{bmatrix}1&0.9&0\\0.9&1&0\\0&0&1\end{bmatrix},
\qquad \bz=(1,1,0.9)^\top,\qquad k=2.
\end{equation}
Diagonal scores select $\{1,2\}$, whose true utility
$\bz_{\cS}^\top\widetilde{\bA}_{\cS}^{-1}\bz_{\cS}$ is approximately $1.0526$.
Support $\{1,3\}$ has utility $1.81$ and therefore lower cost for the same
gain. A block-diagonal or controlled-coherence assumption is necessary.
\end{remark}

\subsection{Expressive Capacity on a Curvature-Weighted Manifold}
\label{app:expressivity}

This subsection makes the ``how much of the update space is occupied, and
relative to which historical demand'' of Section~\ref{sec:expressivity}
precise, and shows that the rank-$R$ capacity statement
(Theorem~\ref{thm:spectral-capacity}) is the rank-$R$ reduction of a
curvature-weighted Grassmann construction. The safe-space contraction
(Theorem~\ref{thm:safe-subspace}) is stated at the ambient kernel and is
\emph{not} carried onto this manifold: with damping $\lambda>0$ the whitening
operator is strictly positive definite, so it has no kernel to identify with the
safe space (Remark~\ref{rem:safe-space-kernel-only}). Throughout, damping
$\lambda>0$ makes every whitening operator invertible so that the embedding is
well defined on all of $\R^p$; the $\lambda\to0$ limit is a separate limiting
statement, not used in the proofs.

\begin{definition}[Curvature-whitened embedding and occupied subspace]
\label{def:whitening-embedding}
Let $\bA_t^{\mathrm{amb}}=\sum_{s\le t}\omega_s\bK_s$ be the accumulated
historical curvature from Eq.~\eqref{eq:additive-safe-curvature}, and let
$\bW_t=(\bA_t^{\mathrm{amb}}+\lambda\bI)^{1/2}\succ0$ for $\lambda>0$. For an
occupied update subspace $\cU\subseteq\R^p$ with orthonormal basis matrix
$\bB_\cU\in\R^{p\times r}$, $\dim(\cU)=r$, the curvature-whitened image is
$\phi_t(\cU)=\range(\bW_t\bB_\cU)\in\mathrm{Gr}(r,p)$, a point of the Grassmann
manifold of $r$-planes in $\R^p$.
\end{definition}

The choice of orthonormal basis $\bB_\cU$ does not affect
$\range(\bW_t\bB_\cU)$, so $\phi_t$ is well defined on subspaces rather than on
bases. A LoRA pool and a full fine-tuning update supply their occupied subspace
$\cU$ differently (architecturally fixed $r=R$ versus optimization-determined
$r=r_{\mathrm{eff}}$); feeding both through $\phi_t$ is an \emph{empirical}
comparison at matched effective rank, not a claim that the rank-$R$ matrix class
(a nonlinear variety) is itself an $r$-plane. We use it only to compare
already-obtained updates, never to identify the two training regimes.

\begin{proposition}[Expressive-equivalence metric]
\label{prop:expressivity-metric}
For $r$-dimensional $\cU,\cV\subseteq\R^p$ with principal angles
$0\le\theta_1\le\cdots\le\theta_r\le\pi/2$ between $\phi_t(\cU)$ and
$\phi_t(\cV)$, define
\begin{equation}
  d_t(\cU,\cV)=\Big(\textstyle\sum_{i=1}^{r}\sin^2\theta_i\Big)^{1/2}.
  \label{eq:app-expressivity-distance}
\end{equation}
Then $d_t$ is a metric on $\mathrm{Gr}(r,p)$, and in particular
$d_t(\cU,\cV)=0$ if and only if $\phi_t(\cU)=\phi_t(\cV)$. Since $\bW_t\succ0$ is
invertible for $\lambda>0$, $\phi_t$ is a linear isomorphism and
$\phi_t(\cU)=\phi_t(\cV)$ holds if and only if $\cU=\cV$; the whitening
therefore does \emph{not} create any new stream-relative collapse of distinct
subspaces at $\lambda>0$. Its role is quantitative, not to enlarge the
equivalence classes: it reweights the principal angles by accumulated historical
curvature, so that for $\cU\neq\cV$ the value $d_t$ orders subspace pairs by
curvature-weighted overlap rather than by raw overlap. Any collapse of distinct
subspaces to $d_t=0$ is a property of the undamped $\lambda\to0$ limit only,
where $\bW_t$ becomes singular on $\ker(\bA_t^{\mathrm{amb}})$; that limit is a
separate statement we do not use in the results below.
\end{proposition}

\begin{proof}
The chordal distance $\big(\sum_i\sin^2\theta_i\big)^{1/2}$ is the standard
Grassmann metric on $\mathrm{Gr}(r,p)$; nonnegativity, symmetry, the triangle
inequality, and the identity of indiscernibles are standard properties of
principal-angle geometry. The principal angles between two $r$-planes with
orthonormal basis matrices $\bQ_{\cU},\bQ_{\cV}$ are
$\arccos\sigma_i(\bQ_{\cU}^\top\bQ_{\cV})$, so $d_t=0$ iff all principal angles are
zero iff $\range(\bQ_{\cU})=\range(\bQ_{\cV})$, which here is
$\range(\bW_t\bB_{\cU})=\range(\bW_t\bB_{\cV})$ since $\bW_t$ is invertible and
preserves dimension; invertibility further forces $\cU=\cV$, giving the stated
identity of indiscernibles at $\lambda>0$. The metric axioms
(nonnegativity, symmetry, triangle inequality) and this identity are therefore
unconditional for every $\lambda>0$. The undamped limit is genuinely different:
as $\lambda\to0$, $\bW_t$ becomes singular on $\ker(\bA_t^{\mathrm{amb}})$, at
which point replacing $\bK_s$ by $c\bK_s$ ($c>0$) leaves the set of directions
with nonzero whitened image unchanged, so a positive rescaling of a single
curvature summand does not move the whitened planes. We record this only as a
property of the limit and do not invoke it in any downstream result, all of
which are stated at fixed $\lambda>0$.
\end{proof}

\begin{remark}[Expressive equivalence is not output-function equivalence]
\label{rem:expressivity-not-function}
$d_t(\cU,\cV)=0$ records that two adapters occupy the same
historically-whitened subspace; it does not constrain the coefficients placed
on that subspace. Two adapters may share $\phi_t(\cU)$ and still realize
different output functions. Expressive equivalence is therefore strictly
weaker than function equivalence and is stated only relative to the
accumulated historical curvature $\bA_t^{\mathrm{amb}}$.
\end{remark}

\begin{proposition}[Effective capacity as a curvature-whitened tail]
\label{prop:effective-capacity}
Let the historical mixture risk satisfy the Kronecker quadratic
Eq.~\eqref{eq:kronecker-quadratic} with local optimum $\bX^\star$, and let
$\cU=\range(\bX-\bX_R^\star)$ collapse to the rank-$R$ class $\cC_R$. Then the
effective capacity in Eq.~\eqref{eq:effective-capacity} satisfies
\begin{equation}
  \kappa_t(\cC_R)
  =\frac12\sum_{j>R}\sigma_j^2(\bY^\star),
  \qquad
  \bY^\star=\bC_{\mathrm{out}}^{1/2}\bX^\star\bC_{\mathrm{in}}^{1/2},
  \label{eq:effective-capacity-reduces}
\end{equation}
which is exactly $\mathsf{Capacity}_{R,T}$ in
Eq.~\eqref{eq:spectral-capacity}. Thus
Proposition~\ref{prop:effective-capacity} reduces to
Theorem~\ref{thm:spectral-capacity} when the occupied subspace is the
rank-$R$ class and the whitening curvature is the Kronecker curvature.
\end{proposition}

\begin{proof}
Under Eq.~\eqref{eq:kronecker-quadratic}, the curvature-whitening operator
of Eq.~\eqref{eq:effective-capacity} is, up to the $\bC_{\mathrm{out}},
\bC_{\mathrm{in}}$ factors, the change of variables
$\bY=\bC_{\mathrm{out}}^{1/2}\bX\bC_{\mathrm{in}}^{1/2}$ used in the proof of
Theorem~\ref{thm:spectral-capacity}. The tail of the local optimum orthogonal
to the rank-$R$ occupied class is, by the Eckart--Young--Mirsky theorem, exactly
the squared singular values of $\bY^\star$ beyond index $R$. Substitution
reproduces Eq.~\eqref{eq:spectral-capacity}, which proves the claimed reduction.
\end{proof}

\begin{remark}[The reduction is to the \emph{optimal} rank-$R$ class, not to an
arbitrary occupied subspace]
\label{rem:effective-capacity-optimal-only}
Proposition~\ref{prop:effective-capacity} is stated and proved only for
$\cU=\cC_R$, the rank-$R$ class that Eckart--Young--Mirsky selects as the best
rank-$R$ approximant. It does \emph{not} generalize to an arbitrary fixed
occupied subspace $\cU$ of dimension $r$. Eckart--Young gives a
\emph{minimum}, not an identity: for a general $\cU$ the whitened residual obeys
\begin{equation}
  \kappa_t(\cU)\;\ge\;\tfrac12\sum_{j>r}\sigma_j^2(\bY^\star),
  \label{eq:effective-capacity-lower-bound}
\end{equation}
with equality \emph{iff} $\cU$ coincides with the top-$r$ left/right singular
subspaces of $\bY^\star$. A concrete gap: with $\bY^\star=\operatorname{diag}(2,1)$
and $\cU=\operatorname{span}(E_{22})$ (a one-dimensional occupied class,
$r=1$, where $E_{22}$ is the matrix unit in coordinate $(2,2)$), projecting
$\bY^\star$ off $\cU$ removes the smaller component and leaves $2E_{11}$, giving
whitened residual $\tfrac12\cdot 2^2=2$, whereas the optimal rank-$1$ tail is
$\tfrac12\sigma_2^2(\bY^\star)=\tfrac12$. Hence $\kappa_t(\cU)=2\neq\tfrac12$; the
tail formula is recovered only at the optimal $\cC_R$. We therefore use Eq.~\eqref{eq:effective-capacity-reduces} exclusively
as the value at the optimal rank-$R$ class and treat
Eq.~\eqref{eq:effective-capacity-lower-bound} as the correct statement for a
generic occupied subspace. The method does not rely on the equality for any
$\cU\neq\cC_R$.
\end{remark}

\begin{remark}[Full fine-tuning is not the $R=\infty$ collapse]
\label{rem:full-ft-not-infinity}
A full fine-tuning update lives in all of $\R^p$ only nominally; its effective
occupied subspace $\cU$ has dimension $r_{\mathrm{eff}}=\dim(\cU)<p$ whenever the
optimization trajectory concentrates energy on fewer than $p$ whitened
directions. Under Definition~\ref{def:whitening-embedding}, full fine-tuning is
therefore located at effective dimension $r=r_{\mathrm{eff}}$ on the same
curvature-weighted Grassmann axis as a rank-$R$ LoRA pool at $r=R$, not at the
$R\to\infty$ limit in which $\kappa_t$ vanishes by construction. We only claim
the placement of the effective \emph{dimension}, not an equality of capacity
values: because full fine-tuning need not occupy the top-$r_{\mathrm{eff}}$
whitened singular subspace, its capacity obeys the lower bound
Eq.~\eqref{eq:effective-capacity-lower-bound} rather than the optimal tail
identity Eq.~\eqref{eq:effective-capacity-reduces} (cf.\
Remark~\ref{rem:effective-capacity-optimal-only}). Whether $r_{\mathrm{eff}}$ is
in fact below $p$ is an empirical property of the stream and the optimizer, not
an assumption of the equivalence above.
\end{remark}

\begin{remark}[Safe-space contraction is stated at the ambient kernel, not via whitening]
\label{rem:safe-space-kernel-only}
One might hope to restate the safe-space contraction of
Theorem~\ref{thm:safe-subspace} as a statement about the curvature-whitened
operator $\bW_t=(\bA_t^{\mathrm{amb}}+\lambda\bI)^{1/2}$, identifying the safe
space with the zero-whitened-norm directions $\{\bm v:\norm{\bW_t\bm v}_2=0\}$.
This identification is only valid in the undamped limit $\lambda\to0$: for every
$\lambda>0$ the operator $\bW_t\succ0$ is strictly positive definite, so its only
zero-norm direction is $\bm v=0$ and the ambient safe space
$\cN_t^{\mathrm{amb}}=\ker(\bA_t^{\mathrm{amb}})$ is \emph{not} recovered from
$\bW_t$. We therefore state the contraction only at the ambient kernel level
(Theorem~\ref{thm:safe-subspace}), where it is exact, and do not carry it onto
the whitened Grassmann manifold, whose embedding $\phi_t$ is defined with
$\lambda>0$ precisely so that it is invertible. The realized $\kappa_t$ is
moreover not asserted monotone in $t$: a newly occupied direction may align with
an existing one and transfer rather than consume capacity, exactly as the
non-cancellation counterexample after
Theorem~\ref{thm:sequential-accumulation} shows for forgetting.
\end{remark}

\subsection{Claim-Status Ledger}

\begin{table}[h]
\caption{Mathematical status of the paper's claims.}
\label{tab:claim-status}
\centering
\small
\begin{tabularx}{\linewidth}{lX}
\toprule
Status & Claims \\
\midrule
Exact & Taylor identity with bounded remainder; three-way decomposition. \\
Exact under an explicit surrogate & Sequential path identity under a
quadratic old loss; spectral tail under Kronecker quadratic curvature;
safe-space contraction under additive PSD curvature. \\
Conditional & Expected accumulation under non-cancellation; top-$k$ optimality
under fixed candidates and block-diagonal curvature; approximation factor under
a spectral cross-block bound. \\
Not claimed & Unconditional monotone forgetting; all forgetting is capacity;
global validity for nonlinear transformers; optimality for unrestricted LoRA
factor training; a solution to general model merging. \\
\bottomrule
\end{tabularx}
\end{table}
